\chapter{dag() --- ASCII DAG Renderer for Bijou}\label{record-426}

\noindent{\footnotesize\color{BookMuted}\path{docs/archive/dag-exploration.md}\par}\smallskip

\section{Vision}\label{record-426-vision}

A pure \texttt{(\allowbreak{}d\allowbreak{}a\allowbreak{}t\allowbreak{}a\allowbreak{})\allowbreak{} \allowbreak{}→\allowbreak{} \allowbreak{}s\allowbreak{}t\allowbreak{}r\allowbreak{}i\allowbreak{}n\allowbreak{}g} component that renders directed acyclic graphs as box-and-arrow ASCII art in the terminal. Auto-layout, edge routing, Unicode box-drawing, theme-aware coloring, graceful degradation.

\begin{CodeBlock}
import { dag } from '@flyingrobots/bijou';

dag([
  { id: 'a', label: 'Theme bridge', edges: ['b', 'c'] },
  { id: 'b', label: 'Static renderers', edges: ['d', 'e'] },
  { id: 'c', label: 'TEA app shell', edges: ['d', 'e', 'f', 'g'] },
  { id: 'd', label: 'RoadmapView', edges: ['h'] },
  { id: 'e', label: 'LineageView', edges: ['h'] },
  { id: 'f', label: 'Input system', edges: ['h'] },
  { id: 'g', label: 'graph.watch()', edges: ['h'] },
  { id: 'h', label: 'Cleanup' },
]);
\end{CodeBlock}

Output:

\begin{CodeBlock}
╭────────────────────╮
│ Theme bridge       │
╰────────────────────╯
          │
          ├───────────────────────╮
          ▼                       ▼
╭────────────────────╮  ╭────────────────────╮
│ Static renderers   │  │ TEA app shell      │
╰────────────────────╯  ╰────────────────────╯
          │                       │
          ├───────────────────────┤
          │           ┌───────────┼───────────┐
          ▼           ▼           ▼           ▼
╭──────────────╮╭──────────────╮╭──────────────╮╭──────────────╮
│ RoadmapView  ││ LineageView  ││ Input system ││ graph.watch()│
╰──────────────╯╰──────────────╯╰──────────────╯╰──────────────╯
          │           │           │           │
          └───────────┴─────┬─────┴───────────┘
                            ▼
                  ╭────────────────────╮
                  │ Cleanup            │
                  ╰────────────────────╯
\end{CodeBlock}

\begin{center}\rule{0.5\linewidth}{0.5pt}\end{center}

\section{API Design}\label{record-426-api-design}

\subsection{Data Types}\label{record-426-data-types}

\begin{CodeBlock}
export interface DagNode {
  id: string;
  label: string;
  edges?: string[];           // IDs of nodes this node points to (children/dependents)
  badge?: string;             // Optional short badge text (e.g., "4h", "DONE", "BLOCKED")
  token?: TokenValue;         // Per-node color override (e.g., green for "ready", red for "blocked")
}

export interface DagOptions {
  nodeToken?: TokenValue;     // Default node border/label color
  edgeToken?: TokenValue;     // Edge/connector color
  highlightPath?: string[];   // Node IDs to highlight (e.g., critical path) — renders with a distinct token
  highlightToken?: TokenValue;// Token for highlighted path
  selectedId?: string;        // Single node drawn with cursor-style highlighting
  selectedToken?: TokenValue; // Override token for selectedId
  nodeWidth?: number;         // Fixed node box width (default: auto from longest label)
  maxWidth?: number;          // Terminal width constraint (default: ctx.runtime.columns)
  direction?: 'down' | 'right'; // Layout direction (default: 'down')
  ctx?: BijouContext;
}

// Accepts SlicedDagSource (adapter-driven) or DagNode[] (small graphs)
export function dag(source: SlicedDagSource, options?: DagOptions): string;
export function dag(nodes: DagNode[], options?: DagOptions): string;
\end{CodeBlock}

\subsection{DagSource Adapter (v0.5.0)}\label{record-426-dagsource-adapter-v0.5.0}

Decouple DAG rendering from any specific in-memory representation. A \texttt{D\allowbreak{}a\allowbreak{}g\allowbreak{}S\allowbreak{}o\allowbreak{}u\allowbreak{}r\allowbreak{}c\allowbreak{}e} represents a potentially unbounded graph --- it has no enumeration method. Use \texttt{d\allowbreak{}a\allowbreak{}g\allowbreak{}S\allowbreak{}l\allowbreak{}i\allowbreak{}c\allowbreak{}e\allowbreak{}(\allowbreak{})} to BFS-walk from a focus node and produce a bounded \texttt{S\allowbreak{}l\allowbreak{}i\allowbreak{}c\allowbreak{}e\allowbreak{}d\allowbreak{}D\allowbreak{}a\allowbreak{}g\allowbreak{}S\allowbreak{}o\allowbreak{}u\allowbreak{}r\allowbreak{}c\allowbreak{}e} that the renderer can consume.

\begin{CodeBlock}
// Unbounded graph adapter — can wrap a database, API, etc.
export interface DagSource {
  has(id: string): boolean;                   // existence check (not enumeration)
  label(id: string): string;                  // display label
  children(id: string): readonly string[];    // outgoing edges
  parents?(id: string): readonly string[];    // incoming edges (required for ancestor traversal)
  badge?(id: string): string | undefined;     // optional badge text
  token?(id: string): TokenValue | undefined; // optional per-node color
}

// Bounded slice — extends DagSource with enumerable IDs
export interface SlicedDagSource extends DagSource {
  ids(): readonly string[];                   // all node IDs in the slice
  ghost(id: string): boolean;                 // is this a boundary placeholder?
  ghostLabel(id: string): string | undefined; // e.g., "... 3 ancestors"
}
\end{CodeBlock}

\textbf{the illustrated sequence below}

% Presentation of original arrow notation. See flow-manifest.json.
\begin{center}
\begin{tikzpicture}[book flow node/.style={draw=BookHeading,line width=.65pt,fill=BookPanel,rounded corners=1.5pt,align=center,inner sep=5pt,minimum height=10mm,font=\fontsize{9}{11}\selectfont,execute at begin node={\hyphenpenalty=10000\exhyphenpenalty=10000}},book flow edge/.style={-{Latex[length=2mm]},draw=BookHeading,line width=.7pt}]
\node[book flow node,text width=96mm] (n0) at (0,0.0) {Source};
\node[book flow node,text width=96mm] (n1) at (0,-1.55) {Slice};
\node[book flow node,text width=96mm] (n2) at (0,-3.1) {Render};
\draw[book flow edge] (n0) -- (n1);
\draw[book flow edge] (n1) -- (n2);
\end{tikzpicture}
\end{center}


% Presentation of original arrow notation. See flow-manifest.json.
\begin{center}
\begin{tikzpicture}[book flow node/.style={draw=BookHeading,line width=.65pt,fill=BookPanel,rounded corners=1.5pt,align=center,inner sep=5pt,minimum height=10mm,font=\fontsize{9}{11}\selectfont,execute at begin node={\hyphenpenalty=10000\exhyphenpenalty=10000}},book flow edge/.style={-{Latex[length=2mm]},draw=BookHeading,line width=.7pt}]
\node[book flow node,text width=96mm] (n0) at (0,0.0) {Dag\allowbreak{}Source: your graph, unbounded};
\node[book flow node,text width=96mm] (n1) at (0,-1.55) {dag\allowbreak{}Slice(): BFS bounded walk, ghost boundaries};
\node[book flow node,text width=96mm] (n2) at (0,-3.1) {dag(): Sugiyama on the bounded result};
\draw[book flow edge] (n0) -- (n1);
\draw[book flow edge] (n1) -- (n2);
\end{tikzpicture}
\end{center}


\begin{CodeBlock}
// External graph (never fully loaded)
const graph: DagSource = {
  has: (id) => db.exists(id),
  label: (id) => db.getTitle(id),
  children: (id) => db.getDeps(id),
  parents: (id) => db.getDependents(id),
};

// BFS touches ~20 nodes, not the whole graph
const slice = dagSlice(graph, 'TASK-42', { depth: 2 });
dag(slice, { ctx });

// Composable: slice a slice
const narrower = dagSlice(slice, 'TASK-42', { depth: 1, direction: 'descendants' });
\end{CodeBlock}

\textbf{\texttt{a\allowbreak{}r\allowbreak{}r\allowbreak{}a\allowbreak{}y\allowbreak{}S\allowbreak{}o\allowbreak{}u\allowbreak{}r\allowbreak{}c\allowbreak{}e\allowbreak{}(\allowbreak{})}} wraps \texttt{D\allowbreak{}a\allowbreak{}g\allowbreak{}N\allowbreak{}o\allowbreak{}d\allowbreak{}e\allowbreak{}[\allowbreak{}]} as a \texttt{S\allowbreak{}l\allowbreak{}i\allowbreak{}c\allowbreak{}e\allowbreak{}d\allowbreak{}D\allowbreak{}a\allowbreak{}g\allowbreak{}S\allowbreak{}o\allowbreak{}u\allowbreak{}r\allowbreak{}c\allowbreak{}e} for backward compatibility:

\begin{CodeBlock}
import { arraySource } from '@flyingrobots/bijou';

const src = arraySource(myNodes);  // SlicedDagSource
dag(src, { ctx });                 // equivalent to dag(myNodes, { ctx })
\end{CodeBlock}

\subsection{Fragment Rendering}\label{record-426-fragment-rendering}

Render a subgraph instead of the full DAG. This is critical for large graphs --- show ancestry of a single node, the critical path, or a neighborhood.

\begin{CodeBlock}
export interface DagSliceOptions {
  direction?: 'ancestors' | 'descendants' | 'both';  // default: 'both'
  depth?: number;                                      // max hops from focus (default: Infinity)
}

// DagSource in → SlicedDagSource out (composable)
export function dagSlice(source: DagSource, focus: string, opts?: DagSliceOptions): SlicedDagSource;

// DagNode[] in → DagNode[] out (backward compatible)
export function dagSlice(nodes: DagNode[], focus: string, opts?: DagSliceOptions): DagNode[];
\end{CodeBlock}

\textbf{Usage:}

\begin{CodeBlock}
// Full graph
dag(allNodes);

// Just the ancestors of "Cleanup"
dag(dagSlice(allNodes, 'cleanup', { direction: 'ancestors' }));

// 2-hop neighborhood around "TEA app shell"
dag(dagSlice(allNodes, 'tea-shell', { direction: 'both', depth: 2 }));

// Critical path highlighted on full graph
dag(allNodes, { highlightPath: ['a', 'b', 'd', 'h'], highlightToken: theme.semantic.error });
\end{CodeBlock}

\textbf{Boundary indicators:} When \texttt{d\allowbreak{}a\allowbreak{}g\allowbreak{}S\allowbreak{}l\allowbreak{}i\allowbreak{}c\allowbreak{}e} cuts edges at the boundary, truncated nodes appear as ghost nodes:

\begin{CodeBlock}
          ╭╌╌╌╌╌╌╌╌╌╌╌╌╌╌╌╌╌╌╌╌╮
          ╎ ... 3 ancestors      ╎    ← dashed border = truncated
          ╰╌╌╌╌╌╌╌╌╌╌╌╌╌╌╌╌╌╌╌╌╯
                    │
                    ▼
          ╭────────────────────╮
          │ TEA app shell      │     ← solid border = in fragment
          ╰────────────────────╯
                    │
          ┌─────────┼─────────┐
          ▼         ▼         ▼
  ╭──────────╮╭──────────╮╭──────────╮
  │ View A   ││ View B   ││ Input    │
  ╰──────────╯╰──────────╯╰──────────╯
\end{CodeBlock}

\begin{center}\rule{0.5\linewidth}{0.5pt}\end{center}

\section{Layout Algorithm: Sugiyama-Lite}\label{record-426-layout-algorithm-sugiyama-lite}

A simplified Sugiyama method optimized for terminal rendering of 5-50 node DAGs. Three phases, \textasciitilde200 lines of TypeScript, zero dependencies.

\subsection{Phase 1: Layer Assignment (Topological BFS)}\label{record-426-phase-1-layer-assignment-topological-bfs}

Assign each node to a horizontal layer (row). Uses longest-path method: push nodes as deep as possible so the graph spreads vertically.

\begin{CodeBlock}
function assignLayers(nodes: DagNode[]): Map<string, number> {
  // 1. Build adjacency + compute in-degree
  // 2. Topological sort (Kahn's algorithm)
  // 3. For each node in topo order:
  //    layer[node] = max(layer[predecessors]) + 1
  //    (roots get layer 0)
}
\end{CodeBlock}

\textbf{Complexity:} O(V + E)

\textbf{Why longest-path over network simplex:} For 5-50 nodes the visual difference is minimal. Longest-path is \textasciitilde20 lines; network simplex is \textasciitilde500. Save complexity budget for edge routing.

\subsection{Phase 2: Column Ordering (Barycenter Heuristic)}\label{record-426-phase-2-column-ordering-barycenter-heuristic}

Within each layer, order nodes to minimize edge crossings. For each node, compute the average column position of its parents (barycenter), then sort by that value.

\begin{CodeBlock}
function orderColumns(layers: string[][], edges: Map<string, string[]>): void {
  // For each layer (top to bottom):
  //   For each node in the layer:
  //     barycenter = average(column indices of parent nodes in previous layer)
  //   Sort layer by barycenter values
  //   (Nodes with no parents keep their original order)

  // Optional: second pass bottom-to-top for better results
}
\end{CodeBlock}

\textbf{Complexity:} O(V log V) per layer per pass. One top-down pass is sufficient for most DAGs. A second bottom-up pass improves crossing minimization at negligible cost.

\subsection{Phase 3: Rendering (Two-Pass Edge Routing)}\label{record-426-phase-3-rendering-two-pass-edge-routing}

The core insight: \textbf{accumulate edge directions first, then render characters.} Never write edge characters during edge processing --- two edges from the same source share vertical exit segments.

\subsubsection{Grid Setup}\label{record-426-grid-setup}

\begin{CodeBlock}
ROW_STRIDE = 6 rows per layer:
  Row 0-2: node box (top border, content, bottom border)
  Row 3:   vertical exit from box
  Row 4:   horizontal elbow turn row
  Row 5:   vertical entry / padding

COL_STRIDE = nodeWidth + 4 (gap between node centers)
\end{CodeBlock}

\textbf{ROW\_STRIDE must be \textgreater= 6.} At RS=5, the elbow turn row collides with the arrowhead destination. RS=6 gives exactly one buffer row. RS=7 adds a padding row (cleaner but wastes vertical space).

\subsubsection{Pass 1: Accumulate Edge Directions}\label{record-426-pass-1-accumulate-edge-directions}

For each edge, mark cells with direction flags (\texttt{U}, \texttt{D}, \texttt{L}, \texttt{R}). Multiple edges sharing a cell accumulate their directions into a \texttt{S\allowbreak{}e\allowbreak{}t\allowbreak{}<\allowbreak{}D\allowbreak{}i\allowbreak{}r\allowbreak{}>}.

\begin{CodeBlock}
type Dir = 'U' | 'D' | 'L' | 'R';
const grid: Set<Dir>[][] = /* 2D array of empty Sets */;

function markEdge(fromCol: number, fromLayer: number, toCol: number, toLayer: number): void {
  const sRow = fromLayer * RS + 3;   // exit below source box
  const dRow = toLayer * RS - 1;     // entry above destination box
  const mid = sRow + 1;              // horizontal turn row

  if (fromCol === toCol) {
    // Straight vertical: mark D/U for each row sRow..dRow-1, then arrowhead at dRow
    for (let r = sRow; r < dRow; r++) markDir(r, fromCol, 'D', 'U');
  } else {
    // Elbow: vertical exit → horizontal turn → vertical entry
    markDir(sRow, fromCol, 'D', 'U');          // exit (may merge with straight edges)

    const [minC, maxC] = fromCol < toCol ? [fromCol, toCol] : [toCol, fromCol];

    // Turn at source column
    markDir(mid, fromCol, fromCol < toCol ? 'R' : 'L', 'U');

    // Horizontal segment
    for (let c = minC + 1; c < maxC; c++) markDir(mid, c, 'L', 'R');

    // Turn at destination column
    markDir(mid, toCol, fromCol < toCol ? 'L' : 'R', 'D');

    // Vertical segment down to destination
    for (let r = mid + 1; r < dRow; r++) markDir(r, toCol, 'D', 'U');
  }
  // Arrowhead at dRow
}
\end{CodeBlock}

\textbf{Why this works for junctions:} When a node has both a straight-down edge and a right-elbow edge, the exit cell accumulates \texttt{\{\allowbreak{}D\allowbreak{},\allowbreak{} \allowbreak{}U\allowbreak{}\}} from the straight edge and \texttt{\{\allowbreak{}R\allowbreak{},\allowbreak{} \allowbreak{}U\allowbreak{}\}} from the elbow. The union \texttt{\{\allowbreak{}D\allowbreak{},\allowbreak{} \allowbreak{}R\allowbreak{},\allowbreak{} \allowbreak{}U\allowbreak{}\}} maps to \texttt{├} --- the correct junction character. This is emergent behavior, not special-cased.

\subsubsection{Junction Character Lookup}\label{record-426-junction-character-lookup}

Convert accumulated direction sets to Unicode box-drawing characters:

\begin{CodeBlock}
const JUNCTION: Record<string, string> = {
  'D': '│',  'U': '│',  'DU': '│',
  'L': '─',  'R': '─',  'LR': '─',
  'DR': '╭', 'DL': '╮', 'RU': '╰', 'LU': '╯',
  'DRU': '├', 'DLU': '┤', 'DLR': '┬', 'LRU': '┴',
  'DLRU': '┼',
};

function junctionChar(dirs: Set<Dir>): string {
  const key = [...dirs].sort().join('');   // canonical: alphabetical (D < L < R < U)
  return JUNCTION[key] ?? '+';
}
\end{CodeBlock}

\subsubsection{Shader-based Serialization (cellAt)}\label{record-426-shader-based-serialization-cellat}

Instead of pre-allocating 2D character/token grids and writing to them in multiple passes, the renderer uses an on-demand \texttt{c\allowbreak{}e\allowbreak{}l\allowbreak{}l\allowbreak{}A\allowbreak{}t\allowbreak{}(\allowbreak{}r\allowbreak{}o\allowbreak{}w\allowbreak{},\allowbreak{} \allowbreak{}c\allowbreak{}o\allowbreak{}l\allowbreak{})} function that computes each cell during serialization. Priority order (highest first):

\begin{enumerate}
\def\labelenumi{\arabic{enumi}.}
\tightlist
\item
  \textbf{Node box} --- look up spatial index by row, check column bounds, return char + token from pre-rendered \texttt{N\allowbreak{}o\allowbreak{}d\allowbreak{}e\allowbreak{}B\allowbreak{}o\allowbreak{}x\allowbreak{}R\allowbreak{}e\allowbreak{}s\allowbreak{}u\allowbreak{}l\allowbreak{}t}
\item
  \textbf{Arrowhead} --- check \texttt{g\allowbreak{}.\allowbreak{}a\allowbreak{}r\allowbreak{}r\allowbreak{}o\allowbreak{}w\allowbreak{}s\allowbreak{}.\allowbreak{}h\allowbreak{}a\allowbreak{}s\allowbreak{}(\allowbreak{}e\allowbreak{}n\allowbreak{}c\allowbreak{}o\allowbreak{}d\allowbreak{}e\allowbreak{}A\allowbreak{}r\allowbreak{}r\allowbreak{}o\allowbreak{}w\allowbreak{}P\allowbreak{}o\allowbreak{}s\allowbreak{}(\allowbreak{}r\allowbreak{}o\allowbreak{}w\allowbreak{},\allowbreak{} \allowbreak{}c\allowbreak{}o\allowbreak{}l\allowbreak{})\allowbreak{})}, return \texttt{▼} + token
\item
  \textbf{Edge} --- check \texttt{g\allowbreak{}.\allowbreak{}d\allowbreak{}i\allowbreak{}r\allowbreak{}s\allowbreak{}[\allowbreak{}r\allowbreak{}o\allowbreak{}w\allowbreak{}]\allowbreak{}?\allowbreak{}.\allowbreak{}[\allowbreak{}c\allowbreak{}o\allowbreak{}l\allowbreak{}]}, return \texttt{j\allowbreak{}u\allowbreak{}n\allowbreak{}c\allowbreak{}t\allowbreak{}i\allowbreak{}o\allowbreak{}n\allowbreak{}C\allowbreak{}h\allowbreak{}a\allowbreak{}r\allowbreak{}(\allowbreak{})} + token
\item
  \textbf{Empty} --- return \texttt{'\allowbreak{} \allowbreak{}'} + null
\end{enumerate}

This avoids allocating O(rows × cols) grid arrays and eliminates separate write passes for edges, arrows, highlights, and nodes.

\begin{center}\rule{0.5\linewidth}{0.5pt}\end{center}

\section{Graceful Degradation}\label{record-426-graceful-degradation}

Three rendering paths, following the bijou component pattern:

\subsection{Interactive Mode (TTY)}\label{record-426-interactive-mode-tty}

Full Unicode box-drawing, ANSI colors via \texttt{c\allowbreak{}t\allowbreak{}x\allowbreak{}.\allowbreak{}s\allowbreak{}t\allowbreak{}y\allowbreak{}l\allowbreak{}e\allowbreak{}.\allowbreak{}s\allowbreak{}t\allowbreak{}y\allowbreak{}l\allowbreak{}e\allowbreak{}d\allowbreak{}(\allowbreak{})}, badges, highlights.

\begin{CodeBlock}
╭────────────────────╮
│ Theme bridge    2h │
╰────────────────────╯
          │
          ├───────────────────────╮
          ▼                       ▼
╭────────────────────╮  ╭────────────────────╮
│ Static renderers   │  │ TEA app shell      │
╰────────────────────╯  ╰────────────────────╯
\end{CodeBlock}

\subsection{Pipe Mode}\label{record-426-pipe-mode}

Indented adjacency list --- machine-parseable, no ANSI, no box drawing:

% Presentation of original arrow notation. See flow-manifest.json.
\begin{center}
\begin{tikzpicture}[book flow node/.style={draw=BookHeading,line width=.65pt,fill=BookPanel,rounded corners=1.5pt,align=center,inner sep=5pt,minimum height=10mm,font=\fontsize{9}{11}\selectfont,execute at begin node={\hyphenpenalty=10000\exhyphenpenalty=10000}},book flow edge/.style={-{Latex[length=2mm]},draw=BookHeading,line width=.7pt}]
\node[book flow node,text width=28mm] (n0) at (0,0) {Theme bridge};
\node[book flow node,text width=28mm] (n1) at (-3,-2) {Static renderers};
\node[book flow node,text width=28mm] (n2) at (3,-2) {TEA app shell};
\node[book flow node,text width=28mm] (n3) at (-5,-5) {Roadmap\allowbreak{}View};
\node[book flow node,text width=28mm] (n4) at (-1.7,-5) {Lineage\allowbreak{}View};
\node[book flow node,text width=28mm] (n5) at (1.7,-5) {Input system};
\node[book flow node,text width=28mm] (n6) at (5,-5) {graph.\allowbreak{}watch()};
\node[book flow node,text width=28mm] (n7) at (0,-8) {Cleanup};
\draw[book flow edge] (n0) -- (n1);
\draw[book flow edge] (n0) -- (n2);
\draw[book flow edge] (n1) -- (n3);
\draw[book flow edge] (n1) -- (n4);
\draw[book flow edge] (n2) -- (n3);
\draw[book flow edge] (n2) -- (n4);
\draw[book flow edge] (n2) -- (n5);
\draw[book flow edge] (n2) -- (n6);
\draw[book flow edge] (n3) -- (n7);
\draw[book flow edge] (n4) -- (n7);
\draw[book flow edge] (n5) -- (n7);
\draw[book flow edge] (n6) -- (n7);
\end{tikzpicture}
\end{center}


\subsection{Accessible Mode}\label{record-426-accessible-mode}

Structured text with depth labels for screen readers:

% Presentation of original arrow notation. See flow-manifest.json.
\begin{center}
\begin{tikzpicture}[book flow node/.style={draw=BookHeading,line width=.65pt,fill=BookPanel,rounded corners=1.5pt,align=center,inner sep=5pt,minimum height=10mm,font=\fontsize{9}{11}\selectfont,execute at begin node={\hyphenpenalty=10000\exhyphenpenalty=10000}},book flow edge/.style={-{Latex[length=2mm]},draw=BookHeading,line width=.7pt}]
\node[book flow node,text width=28mm] (n0) at (0,0) {Layer 1: Theme bridge (2h)};
\node[book flow node,text width=28mm] (n1) at (-3,-2) {Layer 2: Static renderers};
\node[book flow node,text width=28mm] (n2) at (3,-2) {Layer 2: TEA app shell};
\node[book flow node,text width=28mm] (n3) at (-5,-5) {Roadmap\allowbreak{}View};
\node[book flow node,text width=28mm] (n4) at (-1.7,-5) {Lineage\allowbreak{}View};
\node[book flow node,text width=28mm] (n5) at (1.7,-5) {Input system};
\node[book flow node,text width=28mm] (n6) at (5,-5) {graph.\allowbreak{}watch()};
\node[book flow node,text width=28mm] (n7) at (0,-8) {Layer 4: Cleanup (end)};
\draw[book flow edge] (n0) -- (n1);
\draw[book flow edge] (n0) -- (n2);
\draw[book flow edge] (n1) -- (n3);
\draw[book flow edge] (n1) -- (n4);
\draw[book flow edge] (n2) -- (n3);
\draw[book flow edge] (n2) -- (n4);
\draw[book flow edge] (n2) -- (n5);
\draw[book flow edge] (n2) -- (n6);
\draw[book flow edge] (n3) -- (n7);
\draw[book flow edge] (n4) -- (n7);
\draw[book flow edge] (n5) -- (n7);
\draw[book flow edge] (n6) -- (n7);
\end{tikzpicture}
\end{center}
{\scriptsize\color{BookMuted}The source summary says ``8 nodes, 10 edges; the figure follows all 12 listed connections.\par}



\subsection{ASCII Fallback}\label{record-426-ascii-fallback}

When Unicode is unavailable (\texttt{T\allowbreak{}E\allowbreak{}R\allowbreak{}M\allowbreak{}=\allowbreak{}d\allowbreak{}u\allowbreak{}m\allowbreak{}b}, \texttt{N\allowbreak{}O\allowbreak{}\_\allowbreak{}U\allowbreak{}N\allowbreak{}I\allowbreak{}C\allowbreak{}O\allowbreak{}D\allowbreak{}E\allowbreak{}=\allowbreak{}1}), swap the junction table:

\begin{CodeBlock}
const JUNCTION_ASCII: Record<string, string> = {
  'D': '|', 'U': '|', 'DU': '|',
  'L': '-', 'R': '-', 'LR': '-',
  'DR': '+', 'DL': '+', 'RU': '+', 'LU': '+',
  'DRU': '+', 'DLU': '+', 'DLR': '+', 'LRU': '+', 'DLRU': '+',
};
// Arrowhead: 'v' instead of '▼'
// Node boxes: +--+ / |  | instead of ╭──╮ / │  │
\end{CodeBlock}

\begin{center}\rule{0.5\linewidth}{0.5pt}\end{center}

\section{Width Constraints \& Responsive Behavior}\label{record-426-width-constraints-responsive-behavior}

\subsection{Auto-sizing}\label{record-426-auto-sizing}

\begin{CodeBlock}
const nodeWidth = options.nodeWidth ?? Math.max(
  ...nodes.map(n => n.label.length + (n.badge?.length ?? 0) + 4),  // │ label badge │
  16,  // minimum
);
const colStride = nodeWidth + 4;
const maxNodesPerLayer = Math.max(...layers.map(l => l.length));
const totalWidth = maxNodesPerLayer * colStride;
\end{CodeBlock}

\subsection{Terminal Overflow}\label{record-426-terminal-overflow}

When \texttt{t\allowbreak{}o\allowbreak{}t\allowbreak{}a\allowbreak{}l\allowbreak{}W\allowbreak{}i\allowbreak{}d\allowbreak{}t\allowbreak{}h\allowbreak{} \allowbreak{}>\allowbreak{} \allowbreak{}m\allowbreak{}a\allowbreak{}x\allowbreak{}W\allowbreak{}i\allowbreak{}d\allowbreak{}t\allowbreak{}h}:

\begin{enumerate}
\def\labelenumi{\arabic{enumi}.}
\tightlist
\item
  \textbf{Shrink node width} --- truncate labels with \texttt{…}
\item
  \textbf{Reduce gap} --- minimum gap of 2 between nodes
\item
  \textbf{If still too wide} --- switch to vertical (single-column) layout with indentation instead of side-by-side boxes
\end{enumerate}

\subsection{ANSI-Aware Width}\label{record-426-ansi-aware-width}

All width calculations use \textbf{visible length} (strip ANSI escapes before measuring):

\begin{CodeBlock}
function visibleLength(str: string): number {
  return str.replace(/\x1b\[[0-9;]*m/g, '').length;
}
\end{CodeBlock}

Apply colors only when writing content into the grid --- never during layout math.

\begin{center}\rule{0.5\linewidth}{0.5pt}\end{center}

\section{\texorpdfstring{XYPH Integration: \texttt{s\allowbreak{}t\allowbreak{}a\allowbreak{}t\allowbreak{}u\allowbreak{}s\allowbreak{} \allowbreak{}-\allowbreak{}-\allowbreak{}v\allowbreak{}i\allowbreak{}e\allowbreak{}w\allowbreak{} \allowbreak{}d\allowbreak{}e\allowbreak{}p\allowbreak{}s}}{XYPH Integration: }}\label{record-426-xyph-integration-status---view-deps}

With the \texttt{D\allowbreak{}a\allowbreak{}g\allowbreak{}S\allowbreak{}o\allowbreak{}u\allowbreak{}r\allowbreak{}c\allowbreak{}e} adapter, XYPH's quest graph is never duplicated --- bijou reads directly from the quest store:

\begin{CodeBlock}
import { dag, dagSlice, type DagSource } from '@flyingrobots/bijou';

// Wrap the quest store as a DagSource — no DagNode[] copy
const questGraph: DagSource = {
  has: (id) => snapshot.quests.has(id),
  label: (id) => snapshot.quests.get(id)!.title.slice(0, 20),
  children: (id) => snapshot.quests.get(id)!.dependsOn ?? [],
  parents: (id) => snapshot.quests.get(id)!.dependents ?? [],
  badge: (id) => `${snapshot.quests.get(id)!.hours}h`,
  token: (id) => {
    const q = snapshot.quests.get(id)!;
    return q.status === 'DONE' ? theme.semantic.success
         : frontierSet.has(id) ? theme.semantic.primary
         : theme.semantic.muted;
  },
};

// Render a 3-hop neighborhood — never loads the full graph
const slice = dagSlice(questGraph, 'task:BJU-007', { depth: 3, direction: 'ancestors' });
console.log(dag(slice, {
  highlightPath: criticalResult.path,
  highlightToken: theme.semantic.error,
}));
\end{CodeBlock}

\subsection{Reactive DAG in TUI}\label{record-426-reactive-dag-in-tui}

With bijou-tui's TEA runtime + git-warp subscriptions:

\begin{CodeBlock}
const app: App<Model, Msg> = {
  init: () => [model, [watchGraph()]],

  update: (msg, model) => {
    if (msg.type === 'graph-changed') {
      return [{ ...model, snapshot: rebuildSnapshot(msg.diff) }, []];
    }
    return [model, []];
  },

  view: (model) => {
    return flex(
      { direction: 'column', width: model.cols, height: model.rows },
      { basis: 2, content: tabs([...]) },
      { flex: 1, content: (w, h) => viewport({
        width: w, height: h,
        content: dag(buildDagNodes(model.snapshot), {
          highlightPath: model.criticalPath,
          maxWidth: w,
        }),
        scrollY: model.scroll.y,
      })},
      { basis: 2, content: statusLine(model) },
    );
  },
};
\end{CodeBlock}

Add a quest in one terminal → the DAG re-renders in the other with the new node slotted and sorted. That's the dream.

\begin{center}\rule{0.5\linewidth}{0.5pt}\end{center}

\section{Implementation Plan}\label{record-426-implementation-plan}

\subsection{Phase 1: Core Layout Engine (\textasciitilde150 LoC)}\label{record-426-phase-1-core-layout-engine-150-loc}

\begin{itemize}
\tightlist
\item
  \texttt{a\allowbreak{}s\allowbreak{}s\allowbreak{}i\allowbreak{}g\allowbreak{}n\allowbreak{}L\allowbreak{}a\allowbreak{}y\allowbreak{}e\allowbreak{}r\allowbreak{}s\allowbreak{}(\allowbreak{})} --- topological BFS, longest-path
\item
  \texttt{o\allowbreak{}r\allowbreak{}d\allowbreak{}e\allowbreak{}r\allowbreak{}C\allowbreak{}o\allowbreak{}l\allowbreak{}u\allowbreak{}m\allowbreak{}n\allowbreak{}s\allowbreak{}(\allowbreak{})} --- barycenter heuristic (1 top-down + 1 bottom-up pass)
\item
  Grid coordinate mapping (\texttt{l\allowbreak{}a\allowbreak{}y\allowbreak{}e\allowbreak{}r\allowbreak{} \allowbreak{}×\allowbreak{} \allowbreak{}R\allowbreak{}S}, \texttt{c\allowbreak{}o\allowbreak{}l\allowbreak{}u\allowbreak{}m\allowbreak{}n\allowbreak{} \allowbreak{}×\allowbreak{} \allowbreak{}C\allowbreak{}S})
\end{itemize}

\subsection{Phase 2: Edge Router (\textasciitilde100 LoC)}\label{record-426-phase-2-edge-router-100-loc}

\begin{itemize}
\tightlist
\item
  Direction accumulation (\texttt{S\allowbreak{}e\allowbreak{}t\allowbreak{}<\allowbreak{}D\allowbreak{}i\allowbreak{}r\allowbreak{}>} per cell)
\item
  Junction character lookup (15-entry table)
\item
  Elbow routing (see the illustrated routing stages below)
\item
  Arrowhead placement
\end{itemize}

% Presentation of original arrow notation. See flow-manifest.json.
\begin{center}
\begin{tikzpicture}[book flow node/.style={draw=BookHeading,line width=.65pt,fill=BookPanel,rounded corners=1.5pt,align=center,inner sep=5pt,minimum height=10mm,font=\fontsize{9}{11}\selectfont,execute at begin node={\hyphenpenalty=10000\exhyphenpenalty=10000}},book flow edge/.style={-{Latex[length=2mm]},draw=BookHeading,line width=.7pt}]
\node[book flow node,text width=96mm] (n0) at (0,0.0) {vertical exit};
\node[book flow node,text width=96mm] (n1) at (0,-1.55) {horizontal turn};
\node[book flow node,text width=96mm] (n2) at (0,-3.1) {vertical entry};
\draw[book flow edge] (n0) -- (n1);
\draw[book flow edge] (n1) -- (n2);
\end{tikzpicture}
\end{center}


\subsection{Phase 3: Renderer --- Shader-based cellAt (\textasciitilde460 LoC)}\label{record-426-phase-3-renderer-shader-based-cellat-460-loc}

\begin{itemize}
\tightlist
\item
  Node box rendering (label + badge + border) into \texttt{P\allowbreak{}l\allowbreak{}a\allowbreak{}c\allowbreak{}e\allowbreak{}d\allowbreak{}N\allowbreak{}o\allowbreak{}d\allowbreak{}e\allowbreak{}[\allowbreak{}]} with pre-segmented graphemes
\item
  Spatial node index (\texttt{M\allowbreak{}a\allowbreak{}p\allowbreak{}<\allowbreak{}r\allowbreak{}o\allowbreak{}w\allowbreak{},\allowbreak{} \allowbreak{}P\allowbreak{}l\allowbreak{}a\allowbreak{}c\allowbreak{}e\allowbreak{}d\allowbreak{}N\allowbreak{}o\allowbreak{}d\allowbreak{}e\allowbreak{}[\allowbreak{}]\allowbreak{}>}) for O(1) amortized node lookup
\item
  Highlight cell set (\texttt{S\allowbreak{}e\allowbreak{}t\allowbreak{}<\allowbreak{}n\allowbreak{}u\allowbreak{}m\allowbreak{}b\allowbreak{}e\allowbreak{}r\allowbreak{}>} using \texttt{e\allowbreak{}n\allowbreak{}c\allowbreak{}o\allowbreak{}d\allowbreak{}e\allowbreak{}A\allowbreak{}r\allowbreak{}r\allowbreak{}o\allowbreak{}w\allowbreak{}P\allowbreak{}o\allowbreak{}s}) for O(1) edge highlight lookup
\item
  On-demand \texttt{c\allowbreak{}e\allowbreak{}l\allowbreak{}l\allowbreak{}A\allowbreak{}t\allowbreak{}(\allowbreak{}r\allowbreak{}o\allowbreak{}w\allowbreak{},\allowbreak{} \allowbreak{}c\allowbreak{}o\allowbreak{}l\allowbreak{})} query during serialization (priority: node box \textgreater{} arrowhead \textgreater{} edge \textgreater{} empty)
\item
  Run-length grouping serializer: consecutive same-token characters grouped into single \texttt{s\allowbreak{}t\allowbreak{}y\allowbreak{}l\allowbreak{}e\allowbreak{}d\allowbreak{}(\allowbreak{})} calls
\item
  Three output mode paths (interactive / pipe / accessible)
\item
  ASCII fallback junction table
\end{itemize}

\subsection{Phase 4: Fragment Support (\textasciitilde50 LoC)}\label{record-426-phase-4-fragment-support-50-loc}

\begin{itemize}
\tightlist
\item
  \texttt{d\allowbreak{}a\allowbreak{}g\allowbreak{}S\allowbreak{}l\allowbreak{}i\allowbreak{}c\allowbreak{}e\allowbreak{}(\allowbreak{})} --- BFS/DFS from focus node with depth limit
\item
  Ghost node rendering for truncated boundaries (dashed borders)
\item
  Edge count annotation (``\ldots{} 3 ancestors'')
\end{itemize}

\subsection{Phase 5: Width Adaptation (\textasciitilde40 LoC)}\label{record-426-phase-5-width-adaptation-40-loc}

\begin{itemize}
\tightlist
\item
  Auto-sizing from label lengths
\item
  the illustrated sequence below
\item
  \texttt{v\allowbreak{}i\allowbreak{}s\allowbreak{}i\allowbreak{}b\allowbreak{}l\allowbreak{}e\allowbreak{}L\allowbreak{}e\allowbreak{}n\allowbreak{}g\allowbreak{}t\allowbreak{}h\allowbreak{}(\allowbreak{})} for ANSI-safe width math
\end{itemize}

% Presentation of original arrow notation. See flow-manifest.json.
\begin{center}
\begin{tikzpicture}[book flow node/.style={draw=BookHeading,line width=.65pt,fill=BookPanel,rounded corners=1.5pt,align=center,inner sep=5pt,minimum height=10mm,font=\fontsize{9}{11}\selectfont,execute at begin node={\hyphenpenalty=10000\exhyphenpenalty=10000}},book flow edge/.style={-{Latex[length=2mm]},draw=BookHeading,line width=.7pt}]
\node[book flow node,text width=96mm] (n0) at (0,0.0) {Terminal width overflow};
\node[book flow node,text width=96mm] (n1) at (0,-1.55) {label truncation};
\node[book flow node,text width=96mm] (n2) at (0,-3.1) {gap reduction};
\node[book flow node,text width=96mm] (n3) at (0,-4.65) {vertical fallback};
\draw[book flow edge] (n0) -- (n1);
\draw[book flow edge] (n1) -- (n2);
\draw[book flow edge] (n2) -- (n3);
\end{tikzpicture}
\end{center}


\subsection{Total: \textasciitilde1014 LoC + tests}\label{record-426-total-1014-loc-tests}

Tests follow bijou's \texttt{c\allowbreak{}r\allowbreak{}e\allowbreak{}a\allowbreak{}t\allowbreak{}e\allowbreak{}T\allowbreak{}e\allowbreak{}s\allowbreak{}t\allowbreak{}C\allowbreak{}o\allowbreak{}n\allowbreak{}t\allowbreak{}e\allowbreak{}x\allowbreak{}t\allowbreak{}(\allowbreak{})} pattern --- one test per output mode, edge cases (empty, single node, diamond, wide fan-out, skip edges spanning multiple layers, cycles rejected with error).

\begin{center}\rule{0.5\linewidth}{0.5pt}\end{center}

\section{Open Questions}\label{record-426-open-questions}

\begin{enumerate}
\def\labelenumi{\arabic{enumi}.}
\tightlist
\item
  \textbf{Horizontal layout (\texttt{d\allowbreak{}i\allowbreak{}r\allowbreak{}e\allowbreak{}c\allowbreak{}t\allowbreak{}i\allowbreak{}o\allowbreak{}n\allowbreak{}:\allowbreak{} \allowbreak{}'\allowbreak{}r\allowbreak{}i\allowbreak{}g\allowbreak{}h\allowbreak{}t\allowbreak{}'})} --- useful for pipelines. Same algorithm, rotated 90°. Worth including in v1 or defer?
\item
  \textbf{Interactive DAG in bijou-tui} --- cursor navigation over nodes, expand/collapse subtrees, click-to-focus. Separate \texttt{d\allowbreak{}a\allowbreak{}g\allowbreak{}E\allowbreak{}x\allowbreak{}p\allowbreak{}l\allowbreak{}o\allowbreak{}r\allowbreak{}e\allowbreak{}r\allowbreak{}(\allowbreak{})} TEA component or extend \texttt{d\allowbreak{}a\allowbreak{}g\allowbreak{}(\allowbreak{})} with a stateful wrapper?
\item
  \textbf{Maximum practical size} --- at what node count does the ASCII layout become unreadable? Likely \textasciitilde30 nodes in a standard 120-col terminal. Beyond that, should we auto-switch to a condensed format (e.g., collapsed layers)?
\item
  \textbf{Edge label support} --- render edge labels on the horizontal segment? (``depends-on'' is implied in XYPH but other users might want labeled edges.)
\item
  \textbf{Grouping / clusters} --- render a box around a set of nodes (e.g., all tasks in a campaign)? Adds significant layout complexity.
\end{enumerate}

